\part{Teoría de números}
En este capítulo hablaremos de
lo básico que concierne a la
teoría de números de olimpiadas.

No discutiremos demasiado los
problemas que aparecen en omapa,
pero si lo elemental que se aplican
a éstos y varios problemas de prática.

\chapter{Básicos}
\section{Divisibilidad}
La intuición de este concepto
se esclarifica con ejemplos.
Decimos que un entero positivo $a$
es divisible entre otro entero
positivo $n$ si es posible
partir el número $a$ en $n$ partes.
Por ejemplo se nos ocurre que definitivamente
$2$ divide a $4$ ya que es posible
partir el $4$ en $2$ partes iguales.
De manera similar $6$ es divisible
entre $2$ ya que es posible partir
el $6$ en $3 + 3$.
También nos damos cuenta de que
\[6 = 2 + 2 + 2\]
por lo que también es posible partir
al $6$ en $3$ partes iguales, o sea
$3$ divide a $6$.

Por lo que podemos considerar
la siguiente manera de pensar en la
división:

\begin{intuition*}
	Un entero positivo $a$ divide a
	un entero positivo $b$ si es posible
	partir a $b$ en $a$ partes iguales.
\end{intuition*}

Este concepto se puede generalizar
con la misma idea de la siguiente manera:

\begin{definition}
	\label{definition:div}
	Decimos que un entero $a$ divide
	a otro entero $b$ si es posible
	encontrar un entero $z$ tal que
	\[az = b.\]
\end{definition}

\begin{question*}
	¿Cómo se relaciona esto con que
	se posible partir un entero en
	una cantidad entera de partes?
\end{question*}

Notemos también que el número $1$
divide a todos los enteros, porque
es obvio que es posible partir un número
en una sola parte.

A la relación de divisibilidad
se la denota por el símbolo `$\mid$'
y ``$a$ divide a $b$'' se simboliza como
\[a \mid b.\]

De manera similar si $a$ \textbf{no}
divide a $b$ se simboliza como
\[a \nmid b.\]

Otra manera de pensarlo es
con tablas de multiplicar:

\begin{align*}
	            & \vdots \\
	5 \times -3 & = -15  \\
	5 \times -2 & = -10  \\
	5 \times -1 & = -5   \\
	5 \times 0  & = 0    \\
	5 \times 1  & = 5    \\
	5 \times 2  & = 10   \\
	5 \times 3  & = 15   \\
	5 \times 4  & = 20   \\
	            & \vdots
\end{align*}

Podemos decir entonces que
$a \mid b$ si $b$ está en la
tabla de multiplicar de $a$.

\begin{question*}
	¿Cómo se relaciona esto con
	la \autoref{definition:div}?
\end{question*}

\subsection{Propiedades de la división}
Ahora establecemos algunas propiedades
que son fundamentales relacionados a
la división entre dos números enteros.

\begin{proposition}[Transitividad de la divisibilidad]
	\label{proposition:transitivitydiv}
	Si $a \mid b$ y $b \mid c$, entonces $a \mid c$
\end{proposition}

\begin{proposition}[Desigualdad con la divisibilidad]
	Si $a \mid b$, entonces $\abs a \leq \abs b$.
\end{proposition}

Esto tiene sentido ya que
claramente no es posible partir un
entero en más partes que unidades éste
contiene.

\begin{proposition}
	\label{proposition:sumdiv}
	Si $a \mid x$ y $a \mid y$, entonces
	$a \mid x + y$.
\end{proposition}

\begin{proposition}
	Si $a \mid b$, entonces $a \mid zb$
	para todos los enteros $z$.
\end{proposition}

\begin{proposition}
	Si $a \mid b$ si y solo si $a \mid -b$
	si y solo si $-a \mid b$.
\end{proposition}

\subsection*{Problemas para esta sección}
\begin{problem}
Demuestre las últimas cinco
proposiciones (\textit{Nota: } Todas se hacen
con la \autoref{definition:div}).
\end{problem}

\section{Algoritmo de la división}
Este es un teorema que nos permite expresar
cada par de dividendo y divisor $a$ y $b$
como
\[a = bq + r\]
donde $0 \leq r < b$ de manera \textbf{única}.

Una manera de ver por qué esto es cierto
es restringir $a > b > 0$ y ver qué pasa
si tomamos a $a$ como dividendo y $b$
como divisor.
Digamos que $q$ es la mayor cantidad de veces que
podemos sacar $b$ de $a$ sin que este se vuelta
negativo, y al resto que sobra de sacar $q$
veces $b$ de $a$ (que es menor que $b$
ya que de lo contrario podríamos sacar más veces $b$
de $a$)
lo llamamos $r$, y de aquí se ve que
\[a = bq + r.\]

Demostrar formalmente este facto no es muy
complicado y requiere de un pensamiento
muy concreto pero fácil de seguir.

\begin{proposition}[Principio del buen orden]
	\label{proposition:well_order}
	Si $A$ es un subconjunto no vacío
	de los números naturales, entonces $A$
	tiene un mínimo elemento.
\end{proposition}

\begin{question*}
	Convéncete que esto es cierto.
\end{question*}

Dado que queremos encontrar un par $(q, r)$
tal que $a = bq + r$ y $0 \leq r$,
podemos hablar de \textit{todos}
los enteros $q$ que satisfacen que si
$a = bq + r$ entonces $r \geq 0$
y luego tratar de hallar
el menor valor posible de $r = a - bq$.

Por lo que definimos el conjunto
\[A = \left\{a - bq \colon a - bq \geq 0\right\}.\]

Como $A$ es un conjunto de enteros no
negativos, entonces tiene un mínimo
elemento por \autoref{proposition:well_order}
al que llamaremos $r$ y que claramente es
único.

Para probar que $r < b$ podemos ir por
contradicción y suponer que $r \geq b$,
pero esto implica que
\[0 \leq r - b = a - bq - b = a - b(q + 1),\]
y como $q + 1$ es entero, entonces
$a - b(q + 1)$ está en $A$ ya que éste
es mayor o igual que cero,
pero entonces $r - b < r$
está en $A$, por lo que $r$
no era el menor elemento, contradicción.

Por lo que $0 \leq r < b$ y finalizamos.

\section{Números primos}
Hay enteros positivos que no son el
producto de otros dos enteros mayores que $1$,
como por ejemplo el $2$ y el $3$,
vemos que no existen otros enteros
distintos de $1$ que \textit{dividen}
a $2$ a parte de sí mismo.

Esto también pasa con el $5$, revisando número
por número vemos que ninguno divide a este
salvo el $1$ y el $5$. De hecho, hay infinitos
números que cumplen esto, y se los
llama \textit{números primos}.

\begin{definition}
	Un número primo es un entero positivo con
	exactamente dos divisores positivos.
\end{definition}

Como ya vimos, el número $1$ divide a todo número,
y a su vez $n$ siempre divide a $n$, esto implica
que si $n$ es primo, entonces éstos son sus únicos
dos divisores positivos (nótese que así el $1$ no es primo).

\begin{theorem}[Teorema fundamental de la aritmética]
	Todo entero positivo $n \geq 2$ puede ser descompuesto
	en factores primos de manera única.
\end{theorem}

Para desglosar qué dice este teorema,
una manera de verlo es que si $n$ es un
entero positivo y $p_1$, $p_2$, \dots, $p_n$
son todos los primos que lo dividen, entonces
existen unos únicos exponenetes $\alpha_1$, $\alpha_2$,
\dots, $\alpha_n$ tales que
\[n = p_1^{\alpha_1}p_2^{\alpha_2}\cdots p_n^{\alpha_n}.\]

Se puede afirmar con bastante seguridad que
este facto se aplica a casi todos los problemas
de teoría de números.
Este teorema lo probamos en la siguiente sección.

Un buen tip es que muchas veces es mejor
pensar a los números no como
las cantidades que representan, sino
en términos de sus factores primos.

\section{Máximo común divisor}
Existen muchas maneras de definir el
máximo común divisor entre dos
enteros, aquí exhibimos tres de ellas
(que son todas equivalentes).

\begin{definition}
	El máximo común divisor entre dos enteros
	no ambos iguales a cero
	es el producto de los primos comunes (con repetición)
	entre las factorizaciones de éstos.
\end{definition}

\begin{definition}
	El máximo común divisor
	entre dos enteros $a$ y $b$
	no ambos iguales a cero
	es el entero positivo que
	satisface:
	\begin{itemize}
		\ii $d \mid a$ y $d \mid b$.
		\ii Si $c \mid a$ y $c \mid b$, entonces $c \mid d$.
	\end{itemize}
\end{definition}

\begin{definition}
	El máximo común divisor entre dos enteros
	no ambos iguales a cero
	es el mayor entero positivo que
	divide a ambos.
\end{definition}

A este número lo denotamos por $\gcd (a, b)$
y por ser más intuitivo tomaremos
la última definición y probaremos
que es equivalente a las otras dos.

\begin{lemma}[Teorema de Bézout]
	\label{lemma:bezout}
	Si $a$ y $b$ son dos enteros no ambos iguales a cero,
	entonces existen enteros $x$ e $y$ tales que
	\[ax + by = \gcd (a, b).\]
\end{lemma}

Una manera de ver esto es fijarse en la
imagen general y considerar a \textbf{todos}
los enteros de la forma $ax + by$
y probar que $\gcd (a, b)$ debe estar
en ese conjunto. Por lo que lo construimos:
\[S = \left\{ax + by \colon ax + by > 0\right\},\]
y sea $d$ el menor elemento de $S$ (nótese que
$S$ no es vacío). Probaremos que $d = \gcd (a, b)$.
Demostrar directamente esto puede ser difícil,
por lo que en su lugar probaremos que $d \leq \gcd(a, b)$
y que $\gcd (a, b) \leq d$.
Claramente $\gcd (a, b) \mid d$,
por lo que $\gcd (a, b) \leq d$.
\begin{question*}
	¿Por qué pasa esto?
\end{question*}

Ahora solo falta probar que $d \leq \gcd(a, b)$,
y sería suficiente probar que $d \mid a$ y $d \mid b$.
\begin{question*}
	¿Por qué?
\end{question*}
Por construcción existen dos enteros $x_0$ e $y_0$
que satisfacen
\[ax_0 + by_0 = d.\]
Podemos suponer lo contrario, que
$d \nmid a$ o $d \nmid b$, sin pérdida de generalidad
asumimos que $d \nmid a$.
Por el algoritmo de la división
podemos escribir
\[a = dq + r\]
con $0 < r < d$ ya que $d \nmid a$.
Pero esto implica que
\[r = a - dq = a - (ax_0 + by_0) = a(1 - x_0) + (-y_0)b\]
está en el conjunto, que es una contradicción
debido a que $r < d$ y $d$ habíamos supuesto
era el mínimo elemento. Esto implica
que $d \mid a$ y $d \mid b$
completando la prueba.

Aquí un teorema muy útil:
\begin{lemma}[Lema de Euclides]
	\label{lemma:euclidprime}
	Si $p$ es un primo tal que $p \mid ab$,
	entonces $p \mid a$ o $p \mid b$.
\end{lemma}
Tenemos dos casos:
\paragraph{$p \mid a$. } Ya estaríamos.
\paragraph{$p \nmid a$. }
Si esto ocurre, sería suficiente
probar que $p \mid b$.
$p \nmid a$ esto es lo mismo
que decir $\gcd (a, p) = 1$,
ya que caso contrario,
si $d = \gcd(a, p) > 1$, entonces $d = p$
(ya que $p$ es primo y $d \mid p$),
lo que implicaría $p \mid p$ y $p \mid a$
por ser el máximo común divisor,
contradiciendo que $p \nmid a$.

Por el \autoref{lemma:bezout},
existen enteros $x$ e $y$ tales que
\[ax + py = 1,\]
multiplicando ambos lados por $b$
concluimos que
\[abx + pby = b.\]
Como $p \mid ab$, entonces $p \mid b$,
completando la prueba.
\begin{remark}
	Recuerda la \autoref{proposition:sumdiv}.
\end{remark}

\subsection{Demostración del Teorema Fundamental de la Aritmética}
Para probar este teorema debemos
probar cosas:
\begin{itemize}
	\ii Todo entero positivo mayor que
	$1$ puede expresarse como un producto de
	primos.
	\ii Sea cual sea la descomposición
	en factores primos que obtengamos,
	estos serán siempre los mismos primos
	y aparecerán la misma cantidad de veces.
\end{itemize}

El primer ítem requiere de inducción,
siendo la base $n = 2$,
que puede ser expresado como
un producto de números primos ya que este es primo.
Para el paso inductivo supongamos que
todos los enteros positivos $2 \leq n < K$
se pueden expresar como producto de primos,
si $n = K$, hay dos posibilidades:

\paragraph{$n$ es primo. } Si esto ocurre entonces
ya está.

\paragraph{$n$ es compuesto. } Esto significa que
$n$ tiene por lo menos un divisor $a > 0$ que no es $1$ ni $n$,
como $a \mid n$ entonces $a < n$ y existe un
entero $b$ tal que $ab = n$, de la misma manera $b < n$,
por lo que podemos aplicar la hipótesis inductiva a $a$
y $b$ y concluir que éstos pueden ser expresados
como un producto de primos, y de ahí se sigue que $n$
también es un producto de primos, completando
la inducción.

El segundo punto requiere el lema de Euclides.
Supongamos que escribimos a $n$
como un producto de primos de dos maneras:
\[p_1p_2 \cdots p_k = q_1q_2 \cdots q_l.\]

Primero podemos simplificar todos los primos
comunes de tal manera que al final no haya un mismo primo que
esté de ambos lados de la ecuación.
Veremos que en esta situación tendremos
necesariamente $1 = 1$ (todos los primos se
cancelaron).
\begin{question*}
	¿Por qué esto implica que teníamos
	entonces la misma factorización?
\end{question*}

Podemos suponer lo contrario,
que al final tenemos algo como esto:
\[p_1p_2\cdots p_s = q_1q_2\cdots q_t.\]

Obviamente $p_1 \mid q_1 \cdot (q_2\cdots q_t)$,
por lo que $p_1 \mid q_1$ o $p_1 \mid q_2\cdots q_t$
por el \autoref{lemma:euclidprime},
como $q_1 \neq p_1$, entonces
la segunda relación de división es cierta,
por lo que $p_1 \mid q_2q_3 \cdots q_t$,
lo que implica $p_1 \mid q_3\cdots q_t$
y así sucesivamente hasta $p_1 \mid q_t$,
lo que es una contradicción.

\section{Algoritmo de la división Euclidiana}
Este es teorema y método que sirve
para calcular el máximo común divisor
entre dos enteros. Consiste en lo siguiente:

\begin{theorem}
	Si $a = bq + r$ con $0 \leq r < b$,
	entonces $\gcd (a, b) = \gcd (b, r)$.
\end{theorem}

Para ilustrar por qué esto es últil,
básicamene si tomamos $b < a$, entonces
$r < a$, por lo que estaríamos disminuyendo
el valor de una de las variables sí o sí,
lo que nos permite repetir este
proceso hasta que eventualmente uno de los
valores divida al otro, y este primero
será el máximo común divisor entre los primeros
dos números.

\begin{proof}
	Si
	\[a = bq + r\]
	claramente $\gcd (a, b) \mid r$,
	por lo que podemos escribir $a = a^\prime d$,
	$b = b^\prime d$ y $r = r^\prime d$
	siendo $d = \gcd (a, b)$.
	\begin{align*}
		a^\prime \cancel {d} & = b^\prime \cancel{d}q + r^\prime \cancel{d} \\
		a^\prime             & = b^\prime q + r^\prime
	\end{align*}

	Notemos ahora que $a^\prime$
	y $b^\prime$ son coprimos entre sí
	ya que $\gcd(a^\prime, b^\prime) = 1$
	Ahora es suficiente probar que $\gcd (b^\prime, r^\prime) = 1$.
	\begin{question*}
		¿Por qué?
	\end{question*}

	Esto es sencillo ya que en caso contrario,
	si $g = \gcd (b^\prime, r^\prime) > 1$, entonces
	$g \mid a^\prime$ y $g^\prime \mid b^\prime$,
	lo que es una contradicción porque
	$a^\prime$ y $b^\prime$ eran coprimos entre sí.
\end{proof}

Pero nos percatamos de algo:
nunca usamos el dato de que
$0 \leq r < b$; en realidad no era
necesario, por lo que si podemos
descomponer a $a$ en la forma
$a = bq + r$ sean cual sean
los valores de $r$ y $q$,
podemos afirmar con total certeza que
$\gcd(a, b) = \gcd (b, r)$.
Esto deja una consecuencia inmediata:
\begin{corollary}
	$\gcd(a, b) = \gcd(a, b + ak)$, para todo $k \in \ZZ$.
\end{corollary}

Con este teorema en la bolsa, se vuelve
sencillo calcular el máximo común divisor
entre dos números sin la necesidad de factorizarlo.

\begin{example}
	Calcule el máximo común divisor entre
	$124$ y $48$.
\end{example}
Pues simple:
\begin{align*}
	124 & = 48 \cdot 2 + 28 \\
	48  & = 28 \cdot 1 + 20 \\
	28  & = 20 \cdot 1 + 8  \\
	20  & = 8 \cdot 2 + 4   \\
	4   & \mid 8
\end{align*}
Por lo que
\[\gcd (124, 48) = \gcd(28, 48) = \gcd(28, 20) = \gcd(8, 20) = \gcd(8, 4) = 4.\]

A este proceso se lo llama \textit{Algoritmo de Euclides}
para el máximo común divisor, y es de lo que trata
esta sección. Más específicamente el algoritmo escrito
en forma general para calcular el máximo
común divisor entre $r_0$ y $r_1$ ($\abs{r_0} \geq \abs{r_1}$) es así:

% ISSUE: The \vdots looks off-center with respect 
% to the equal signs
\begin{minipage}{\textwidth}
	\begin{align*}
		r_0 & =      r_1q_0 + r_2      \\
		r_1 & =      r_2q_1 + r_3      \\
		    & \vdots                   \\
		r_k & =      r_{k + 1}q_k + 0.
	\end{align*}
	siendo $\gcd(r_0, r_1) = r_{k + 1}$.
\end{minipage}

\chapter{Aritmética modular}
Esta es una idea en la que se sostiene
casi todo lo desarrollado respecto
a la teoría de números, y consiste en
categorizar qué en general
si analizamos los restos que deja
el algoritmo de división euclidiana.

Para nosotros este pedazo de teoría
nos ayudará a abreviar muchas cosas.

El clásico ejemplo para describir a la
aritmética modular es el reloj de $12$ horas:

% TODO: Insertar imagen de reloj de 12 horas

Digamos que el reloj empieza en la hora $0$,
entonces al cabo de $10$ horas el reloj
marcará, $10$, al cabo de $11$ marcará $11$,
pero al paso de $12$ horas éste volverá a marcar $0$
\footnote{Esquivada olímpica.}. % TODO: Quitar este footnote.

Esto se puede formalizar a través de
analizar los restos en \textit{módulo} $12$
y decimos que las horas que marca el reloj
son los enteros \textit{módulo} $12$.

\begin{definition}
	Dos enteros son congruentes módulo $n$
	si tienen el mismo resto divididos entre $n$.
\end{definition}

Esto abre paso a, quizás, querer agrupar
a los enteros teniendo en cuenta
su resto al dividirse entre $n$,
en particular, no es difícil de ver que:

\begin{theorem}\label{theorem:congruence}
	Dos enteros $a$ y $b$ son congruentes
	módulo $n$ si y solo si $n$ divide a $a - b$.
\end{theorem}

Más aún, ver a los enteros desde esta perspectiva
nos abre un gran abanico de herramientas
para entenderlos, entre ellos se dejan ciertas
propiedades como ejercicios al lector.

\begin{definition}
	Denotamos a `$\equiv$' como el símbolo
	de congruencia. La expresión
	``$a$ y $b$ son congruentes módulo $n$''
	se escribe matemáticamente de la siguiente
	manera:
	\[a \equiv b \pmod n\]
\end{definition}

\subsection*{Problemas para esta sección. }
\begin{problem}
Demuestre el \autoref{theorem:congruence}.
\end{problem}

\begin{problem}[Transitividad de la congruencia]
Demuestre que si $a \equiv b$ y $b \equiv c$,
entonces $a \equiv c$.
\end{problem}

\begin{problem}[Suma de congruencias]
Demuestre que si $a \equiv b$ y $x \equiv y$ en
módulo $n$, entonces
\[a + x \equiv b + y \pmod n.\]
\end{problem}

\begin{problem}[Congruencia y multiplicaciones]
\label{theorem:congprod}
Demuestre que si $a \equiv x$ y $b \equiv y$
en módulo $n$ entonces
\[ab \equiv xy \pmod n.\]
\end{problem}

\begin{problem}[Congruencia y la exponenciación]
Demuestre que si $a \equiv b$ en módulo $n$,
entonces
\[a^k \equiv b^k \pmod n\]
para todo entero positivo $k$.

\begin{hints}
	\begin{hint}
		\autoref{theorem:congprod}.
	\end{hint}
\end{hints}
\end{problem}

\chapter{Problemas}
\begin{problem}
Alicia y Beto juegan a un juego
por turnos de forma alternada.
En cada turno el jugador
escribe un dígito del $1$
al $9$ en la pizarra.
Si ambos juegan la misma cantidad de veces
y empieza Alicia,
¿puede ella asegurar que después
de escribir todos los dígitos
sea posible escribir un número primo
con esos dígitos?

\begin{hints}
	\begin{hint}
		¿Puede Beto asegurar que el número sea
		múltiplo de algo?
	\end{hint}

	\begin{hint}
		Múltiplos de $3$.
	\end{hint}
\end{hints}
\end{problem}

\begin{problem}
Calcule la cantidad de divisores enteros positivos de $2026$.
\end{problem}

% \begin{problem}
% Calcule la suma de los divisores naturales de $2026$.
% \end{problem}

\begin{problem}
Encuentre la cantidad de triplas de enteros positivos $a$, $b$ y $c$ tales que $a$ divide a $bc + 1$, $b$ divide
a $ac + 2$ y $c$ divide a $ab + 3$.
\end{problem}

\begin{problem}
Encuentre el resto de dividir $2026^{2026^{2026}}$ entre $7$.
\end{problem}

\begin{problem}
Encuentre todas las ternas $a$, $b$, $c$ de enteros positivos
tales que
\[a^2 + b^2 + 1 = c!.\]
\end{problem}

\begin{problem}
Demuestre que que para todo $n$ natural se cumple que
\[6 \mid n^3 + 5n.\]
\end{problem}

\begin{problem}
Encuentre todos los pares de enteros $\left(x, y\right)$ tales que
\[\frac{xy}{x + y} = 144\]
\end{problem}
%
\begin{problem}[OMAPA 2025/1]
Llamamos a un número \textit{lomaplatense} si tiene tres dígitos y al sumarlo
con el número que resulta de invertir orden de sus dígitos,
el resultado es otro número de tres dígitos con todos sus dígitos
iguales. Por ejemplo, $210$ es lomaplatense ya que $210 + 012 + 222$.
¿Cuántos números \textit{lomaplatenses} hay?
\begin{hints}
	\begin{hint}
		Consiste en analizar la suma $\overline{abc} + \overline{cba}$.
	\end{hint}

	\begin{hint}
		$a + c \leq 9$.
	\end{hint}
\end{hints}
\solutionlabel{problem:omapa2025-1}
\end{problem}

\begin{problem}[Selectivo IMO 2026/1]
Encuentre todas las ternas de enteros no negativos
$(p, x, y)$ con $p$ primo que satisfacen
\[p^x = y^4 + 4.\]
\begin{hints}
	\begin{hint}
		Completa el cuadrado.
	\end{hint}
	\begin{hint}
		Diferencia de cuadrados.
	\end{hint}
	\begin{hint}
		Lema de Euclides.
	\end{hint}
\end{hints}
\solutionlabel{problem:selectivo20261}
\end{problem}

\begin{solutiontoproblem}{problem:selectivo20261}
	Notemos que
	\[y^4 + 4 = y^4 + 4y^2 + 4 - 4y^2 = (y^2 + 2)^2 - (2y)^2 = (y^2 + 2y + 2)(y^2 - 2y + 2) = p^x.\]
	Claramente $y^2 + 2y + 2 > 1$.
	\paragraph{Caso 1. $y^2 - 2y + 2 = 1$.}
	Esto implica $y = 1$, y sustituyendo nos queda que necesariamente $p = 5$ y $x = 1$
	es una solución.
	\paragraph{Caso 2. $y^2 - 2y + 2 > 1$. }
	Esto implica que $p \mid y^2 - 2y + 2$ y $p \mid y^2 + 2y + 2$,
	por lo que
	\[p \mid y^2 + 2y + 2 - (y^2 - 2y + 2) = 4y.\]
	Por el lema de Euclides, ocurre una de dos cosas:
	\paragraph{Caso 2.1. $p \mid 4$. } Esto implica
	que $p = 2$, lo que implica que $y$ es par,
	pero si $y \geq 2$, entonces $16 \mid y^4$
	por lo que
	\[2^x = y^4 + 4 = 4(4t + 1)\]
	que es imposible ya que $4t + 1$ no puede
	ser potencia de $2$ con $t \geq 1$.
	Esto implica que $y = 0$, por lo que $x = 2$.

	\paragraph{Caso 2.2. $p \mid y$. } En la ecuación original
	vemos que entonces $p \mid 4$ y regresamos al \textbf{caso 2.1}.
	%
	En síntesis, las únicas dos ternas soluciones son
	\[(5, 1, 1) \  \text{y} \  (2, 2, 0).\]
\end{solutiontoproblem}

\begin{problem}[CAM JT 2025]
Demuestre que si $\overline{abc\dots{x}yz}$ es múltiplo de $37$,
entonces $\overline{z0y0x0\dots0c0b0a}$ es múltiplo de 37.
(Ej.: $259$ es múltiplo de $37$, entonces $90502$ es múltiplo de 37).

%%TODO: Poner el link de referencia 
\textit{Nota: } Encontré este problema \href{https://www.youtube.com/watch?v=d6iQrh2TK98&t=178s}{aquí}.
\begin{hints}
	\begin{hint}
		$10^3 \equiv 1 \pmod {37}$.
	\end{hint}
	\begin{hint}
		Expresa $10^na_n + \cdots + 10a_1 + a_0 = \overline{a_n\dots a_1a_0} \equiv 0 \pmod{37}$.
	\end{hint}
	\begin{hint}
		Analice el patrón que generan las potencias
		de $10$ módulo $37$.
	\end{hint}
\end{hints}
\end{problem}

\begin{problem}[OMAPA 2011/3]
Se divide $\ol{aaaa}$ entre $\ol{bb}$ obteniéndose
un cociente entre $140$ y $160$ inclusive y como residuo
$\ol{(a - b)(a - b)}$. Determine todos los pares $(a, b)$
que cumplen esta condición.
\end{problem}

\begin{problem}
Encuentre todas las funciones $f\colon\NN\to\NN$ tales que $f(2) = 2$, $f$
es estrictamente creciente y $f(xy) = f(x)f(y)$ para todos los pares de
enteros positivos $(x, y)$.
\end{problem}

\begin{problem}
Sea $(f_n)$ la sucesión de Fibonacci.
Demuestre que cualesquiera dos
términos consecutivos son coprimos entre sí.
\end{problem}

\begin{solutiontoproblem}{problem:omapa2025-1}
	Consideremos un número lomaplatense expresado de la siguiente manera: $\overline{abc}$.
	Debemos analizar la suma $\overline{abc} + \overline{cba}$.

	\paragraph{Analisis cuando $a + c \geq 10$.} Esto es imposible
	porque implica que $\overline{abc} + \overline{cba} \geq 1000$
	(suma de los dígitos de las centenas).

	\paragraph{Análisis cuando $2b \geq 10$.} Esto también es imposible
	porque considerando $a + c \leq 9$ por el caso anterior,
	esto implicaría un acarreo al dígito
	de las centenas, haciendo que este sea distinto al dígito de
	la unidad.

	Llegamos a condiciones necesarias y suficientes que debe
	cumplir un número lomaplatense, ($b$ queda totalmente determinado
	por $a$ y $c$):
	\begin{itemize}
		\ii $a + c$ es par.
		\ii $a + c \leq 8$.
		\ii $a \geq 1$.
	\end{itemize}

	La cantidad de pares $\left(a, c\right)$
	que satisfacen $a + c = 2x$ es
	\[2x + 1\]
	pero hace falta restar $1$ porque $a = 0$
	no es posible (corresponde a un solo par).
	Queda sumar para cada $1 \leq b \leq 4$:
	\[2(1 + 2 + 3 + 4) = 20.\]
\end{solutiontoproblem}

\begin{problem}[OMAPA 2010/4]
Encuentre todos los números naturales de cuatro cifras $\ol{abcd}$ que
sean múltiplos de $11$, el número de dos cifras $\ol{ac}$ sea
múltiplo de $7$ y $a + b + c + d = d^2$.
\end{problem}

\begin{problem}[OMAPA Ronda departamental 2011/6]
Sean $M$ y $N$ dos números enteros positivos
con $M^2 - N^2 = 2011$,
¿cuál es el máximo valor que puede tener $M - N$?
\end{problem}
